\documentclass[11pt]{article}

\usepackage[margin=1in]{geometry}
\usepackage{booktabs}
\usepackage{longtable}
\usepackage{array}
\usepackage{enumitem}
\usepackage{float}
\usepackage{graphicx}
\usepackage{amsmath}
\usepackage[hidelinks,pdfborder={0 0 0}]{hyperref}

\newcommand{\code}[1]{\texttt{\detokenize{#1}}}
\newcommand{\pct}[1]{#1\%}
\setlength{\emergencystretch}{3em}

\title{Monthly Milestone Report -- May 2026}
\author{Zihao Wang}
\date{}

\begin{document}
\maketitle

\noindent\begin{tabular}{@{}>{\bfseries}p{1.35in}p{4.75in}@{}}
Research Area: & WiFi CSI human activity recognition, PA-CSI domain generalization, contrastive representation learning, and cross-environment failure analysis \\
Main Focus This Month: & Testing prototype/margin/alignment methods for the $E2$ class $1/4$ failure, separating smoke signals from matched full results, and deciding which directions should be stopped \\
\end{tabular}

\paragraph{Monthly Summary.}
This month changed the project from method expansion to evidence-based pruning. At the beginning of May, the main hypothesis was that the $E2$ failure could be repaired by better contrastive geometry: sub-center prototypes, fixed or adaptive PairMargin, and source-calibrated hard negative selection. By the end of the month, matched multi-seed evidence showed that this explanation is not sufficient. $K=3$ is not significantly better than $K=1$ under current code, geometry-balance variants collapse their active occupancy, K=1 PairMargin hurts the class-balance metric, OT alignment compresses the target $1/4$ geometry, and the CC-VREx smoke gain fails under full validation. The strongest supported conclusion is that the current $E2$ class $1/4$ problem is a target representation collapse problem, not a simple source-side prototype-capacity or margin-tuning problem.

\clearpage

\section{Review of Previously Defined Goals and Progress}

\begin{longtable}{@{}p{1.65in}p{1.05in}p{3.4in}@{}}
\toprule
Previous Goal & Status & Progress Summary \\
\midrule
\endhead
Run the selected \code{K=3 + PairMargin} prototype configuration on full LOEO targets $E1$, $E2$, and $E3$ & Reframed & The original full-LOEO promotion was stopped because the current-code evidence no longer supported multi-center geometry as the main explanation. A matched $K=3$ versus $K=1$ check on $E2$ seeds 42/52/62 showed no significant gain from $K=3$. \\
Compare the full-LOEO prototype result against the old SupCon reference & Reframed & The comparison target remained important, but May work showed that single-target $E2$ gains are not enough. The clean SupCon line is now treated as the anchor, while all new methods must first beat their own matched baselines. \\
If the prototype result is mixed, separate sub-center effects from PairMargin effects & Completed with negative outcome & I ran or aggregated focused checks for $K=1$, $K=3$, geometry balance, K=1 PairMargin, dynamic/adaptive margins, OT, and CC-VREx. The result is a clear negative decision for several added mechanisms. \\
Analyze asymmetric $1\leftrightarrow4$ behavior on $E2$ & Completed & The diagnostics show that the direction of $1\leftrightarrow4$ confusion changes across seeds, and many target mistakes are high-confidence and closer to the predicted source centroid. This argues against hand-built asymmetric margin grids. \\
Convert April results into manuscript-ready tables and figures & Partially completed & I produced May interim tables, experiment summaries, method-result inventories, and an explicit abandonment report. The material is better organized, but it is not yet manuscript-ready because the main contribution boundary changed. \\
Start robustness experiments after the prototype result is known & Not completed & Robustness experiments were deferred. The more urgent issue is now to determine whether the $E2$ class $1/4$ collapse is data/label-related or representation-level domain shift. \\
\bottomrule
\end{longtable}

\paragraph{Overall assessment.}
The month did not deliver a new winning method, but it delivered an important research correction. The project now has a cleaner boundary between supported results, smoke-only signals, and directions that should be stopped. This is useful because it prevents the next month from spending compute on mechanisms that are already contradicted by matched evidence.

\section{Papers and Method Families Reviewed During the Month}

\begin{itemize}[leftmargin=*]
    \item \textbf{Sub-center prototype and angular-margin methods.}
    I continued using sub-center and margin-based recognition ideas as the motivation for prototype capacity and local pair margins. The May result shows that these ideas are not enough for the current $E2$ collapse when applied only as source-side geometry repair.

    \item \textbf{V-REx / risk extrapolation ideas.}
    I tested whether class-conditional source risk disagreement on classes 1 and 4 could improve target robustness. The smoke result was positive, but the matched full run showed that the current CC-VREx configuration is unstable.

    \item \textbf{Optimal transport and target alignment.}
    I evaluated whether unsupervised OT alignment could repair the target $1/4$ collapse without using target labels for training or selection. The full result was negative, and the target $1/4$ centroid distance became smaller rather than larger.

    \item \textbf{MMD/CORAL-style alignment as follow-up candidates.}
    These remain possible alignment baselines, but May OT results suggest that they should not be run as automatic sensitivity experiments. A new representation or pseudo-label mechanism is needed first.

    \item \textbf{PA-CSI and SupCon as anchor methods.}
    The clean PA-CSI + SupCon result remains the most useful anchor. May work reinforced that clean class-aware SupCon is not the problem; the failed add-ons are the problem.
\end{itemize}

\section{What I Learned}

\begin{enumerate}[leftmargin=*]
    \item \textbf{Matched baselines matter more than isolated numbers.} Several May methods look good when compared to a weak or unrelated baseline, but fail when compared against their own matched control.

    \item \textbf{Smoke runs are not evidence of success.} K=1 PairMargin and CC-VREx both had useful smoke signals. Later checks showed that the PairMargin smoke hurt class balance and the CC-VREx full run failed across three seeds.

    \item \textbf{The $E2$ class $1/4$ issue is target-side collapse.} Source class 1 and class 4 are well separated, while target class 1 and class 4 are compressed into a shared region. This is stronger evidence than any single method result.

    \item \textbf{More prototypes do not solve the current failure.} $K=3$ is not significantly better than $K=1$, and geometry-balanced multi-center variants lose active occupancy during training.

    \item \textbf{Manual pair margins can make the wrong trade-off.} The K=1 PairMargin smoke improved accuracy and macro-F1 but reduced $\min(\mathrm{F1}_1,\mathrm{F1}_4)$ from 0.5368 to 0.3654.

    \item \textbf{Ordinary OT alignment is not enough.} OT did not improve the target $1/4$ geometry. It compressed the final target centroids and reduced the focus-pair F1 metric on average.

    \item \textbf{The next step is diagnosis, not another small regularizer.} The project should now ask whether the collapse is caused by data/semantic noise or representation-level domain shift.
\end{enumerate}

\section{Main Experimental Results}

\subsection{Current Reference Lines}

\begin{table}[H]
\centering
\small
\caption{Clean reference lines for interpreting May experiments on target $E2$.}
\label{tab:anchors}
\resizebox{\textwidth}{!}{%
\begin{tabular}{llcccl}
\toprule
Anchor & Scope & Seeds & Acc. & Macro-F1 & Note \\
\midrule
\code{baseline_pa_csi_dg_paper_aligned} & CE baseline, $E2$ & 42 & 0.8547 & 0.8304 & clean single-seed CE anchor \\
\code{pa_csi_dg_supcon_paper_aligned} & SupCon, $E2$ & 42 & 0.8953 & 0.8763 & cleanest single-seed SupCon anchor \\
\code{control_new_kall1} & K=1 source-only control, $E2$ & 42/52/62 & 0.8821 & 0.8582 & matched control for geometry-balance cycle \\
\code{e2_14_phase1_k1_source_only} & K=1 source-only control, $E2$ & 42/52/62 & 0.8508 & 0.8302 & matched control for OT cycle \\
\bottomrule
\end{tabular}
}%
\end{table}

\subsection{May Method Inventory}

\begin{table}[H]
\centering
\small
\caption{May method results and decisions.}
\label{tab:method_inventory}
\resizebox{\textwidth}{!}{%
\begin{tabular}{p{1.65in}p{1.25in}p{2.75in}p{1.2in}}
\toprule
Method & Evidence level & Key result & Decision \\
\midrule
Sub-center prototype + PairMargin & Single-target / historical positive & $K=3$ + PairMargin reached 0.9023 accuracy and 0.8824 macro-F1 on $E2$ seed 42, but later matched checks did not support multi-center as a robust explanation. & Do not promote as main line \\
$K=3$ vs. $K=1$ paired check & Matched 3-seed diagnostic & Accuracy delta $K3-K1=-0.0033$, macro-F1 delta $=-0.0030$, both with $p>0.7$. & Reject K>1 premise \\
Geometry balance / diversity & Matched 3-seed operational gate & \code{geom_balance} and \code{geom_balance_div} aborted at epoch 100 for all seeds due active occupancy below 0.05. & Stop \\
K=1 PairMargin & Smoke / direction check & Accuracy rose, but $\min(\mathrm{F1}_1,\mathrm{F1}_4)$ fell from 0.5368 to 0.3654 and class 1 collapsed further. & Stop this form \\
Dynamic margin / LCA-KM & Single-seed / exploratory & Dynamic and adaptive variants ran, but did not beat the clean SupCon anchor on focus-pair balance. & Keep as future idea only \\
OT alignment & Matched 3-seed full run & Baseline focus-pair F1 mean 0.5888; unfiltered OT 0.5433; prefiltered OT 0.5504. Target $1/4$ geometry compressed. & Stop ordinary OT line \\
CC-VREx focus 1/4 & Matched 3-seed full run & Baseline accuracy/macro-F1/focus-F1: 0.8779/0.8529/0.5996. CC-VREx: 0.8339/0.8001/0.5530. & Stop this configuration \\
\bottomrule
\end{tabular}
}%
\end{table}

\subsection{Matched Negative Results}

\begin{table}[H]
\centering
\small
\caption{Matched full or matched diagnostic results that changed the May decision.}
\label{tab:matched_negative}
\begin{tabular}{lrrr}
\toprule
Comparison & Metric & Baseline / control & Proposed / delta \\
\midrule
$K=3$ vs. $K=1$ & Accuracy delta & -- & -0.0033 \\
$K=3$ vs. $K=1$ & Macro-F1 delta & -- & -0.0030 \\
Unfiltered OT vs. baseline & $\min(\mathrm{F1}_1,\mathrm{F1}_4)$ & 0.5888 & 0.5433 \\
Prefiltered OT vs. baseline & $\min(\mathrm{F1}_1,\mathrm{F1}_4)$ & 0.5888 & 0.5504 \\
CC-VREx vs. matched baseline & Accuracy & 0.8779 & 0.8339 \\
CC-VREx vs. matched baseline & Macro-F1 & 0.8529 & 0.8001 \\
CC-VREx vs. matched baseline & $\min(\mathrm{F1}_1,\mathrm{F1}_4)$ & 0.5996 & 0.5530 \\
\bottomrule
\end{tabular}
\end{table}

\subsection{Representation Collapse Evidence}

\begin{table}[H]
\centering
\small
\caption{Centroid distance evidence for target $E2$ class $1/4$ collapse under the $K=1$ diagnostic line.}
\label{tab:collapse}
\begin{tabular}{rrrr}
\toprule
Seed & Source $1\leftrightarrow4$ distance & Target $1\leftrightarrow4$ distance & Target/source ratio \\
\midrule
42 & 1.2177 & 0.2722 & 0.224 \\
52 & 1.2045 & 0.1528 & 0.127 \\
62 & 1.2540 & 0.2076 & 0.166 \\
\bottomrule
\end{tabular}
\end{table}

The collapse evidence is the most important May result. It explains why source-side prototype and margin methods keep failing. The source representation separates class 1 and class 4, but target $E2$ maps them into a much tighter shared region. A method that only changes source geometry cannot reliably unfold a collapsed target representation.

\section{Identified Problems / Questions}

\begin{itemize}[leftmargin=*]
    \item The current evidence cannot distinguish label/semantic noise from severe target-domain covariate shift. Both would produce high-confidence target mismatches.
    \item The manual-review step is now a blocker. I need to inspect high-confidence unanimous mismatch samples before running another large training sweep.
    \item The success metric should include geometry, not only F1. A future method should show that target class $1/4$ centroid distance increases before it is treated as a real fix.
    \item The project needs a clearer artifact boundary: clean anchors, stopped methods, smoke-only probes, and future candidates should not be mixed in one table.
    \item The final contribution may shift from ``new prototype method'' to ``protocol-aware diagnosis of PA-CSI DG failure plus a more justified representation/alignment method.''
    \item Cross-dataset validation should become the June priority. A new dataset such as StanFi can test whether the current conclusions are specific to PA-CSI MultiEnv $E2$ or reflect a broader domain-generalization issue.
\end{itemize}

\section{Potential Research Directions}

\begin{itemize}[leftmargin=*]
    \item \textbf{Manual target audit:} inspect high-confidence unanimous $1\rightarrow4$ and $4\rightarrow1$ mismatch samples to decide whether the labels/semantics are questionable.
    \item \textbf{Representation-level domain-shift diagnosis:} track raw-space, fused-feature, projection-space, and classifier-logit separation for classes 1 and 4.
    \item \textbf{Pseudo-label construction with confidence and geometry checks:} if the issue is representation shift, design a pseudo-label mechanism that first improves target $1/4$ geometry.
    \item \textbf{Cross-dataset validation with StanFi or a similar CSI HAR dataset:} use a new dataset to check whether the clean SupCon anchor, stopped methods, and collapse diagnostics transfer beyond the current PA-CSI MultiEnv split.
    \item \textbf{Source-only adaptive calibration, redesigned:} keep LCA-KM-style calibration only if it is rebuilt around the collapse evidence instead of reused as another margin grid.
    \item \textbf{Negative-result manuscript material:} the May results can become a useful section on why several intuitive DG fixes fail under matched baselines.
    \item \textbf{Clean SupCon as the stable baseline:} continue treating PA-CSI + SupCon as the anchor unless a matched rerun disproves it.
\end{itemize}

\clearpage
\section{Goals for Next Month}

\begin{longtable}{@{}cp{5.75in}@{}}
\toprule
\# & Goal \\
\midrule
\endhead
1 & Manually review high-confidence unanimous $E2$ class $1/4$ mismatch samples, balanced between $1\rightarrow4$ and $4\rightarrow1$, so the PA-CSI MultiEnv failure has a clear interpretation before external validation. \\
2 & Select and prepare a new CSI HAR validation dataset, with StanFi as the first candidate if the data format and labels can be mapped cleanly to the current pipeline. \\
3 & Implement a dataset adapter and preprocessing bridge for the new dataset: sample tensor shape, amplitude/phase handling, labels, subject/environment split metadata, and train/validation/test manifests. \\
4 & Run the clean baseline and clean SupCon anchor on the new dataset before testing any May add-on method. The first goal is to establish whether the source-only SupCon advantage transfers. \\
5 & Compare PA-CSI MultiEnv and the new dataset using the same reporting discipline: matched baselines, class-wise F1, confusion matrices, and geometry diagnostics. \\
6 & Freeze the stopped-method list for PA-CSI MultiEnv: multi-center repair, occupancy-floor tuning, K1 PairMargin grids, ordinary OT sensitivity, and the current CC-VREx configuration should not be resumed without a new mechanism reason. \\
7 & Rebuild the report tables so that every future result is grouped by dataset, clean anchor, matched full validation, smoke-only probe, or stopped/debug artifact. \\
\bottomrule
\end{longtable}

\clearpage
\section{Six-Month Research and Publication Plan}

\begin{longtable}{@{}p{0.95in}p{5.35in}@{}}
\toprule
Period & Plan \\
\midrule
\endhead
June 2026 & Complete the $E2$ class $1/4$ manual audit, freeze stopped May methods, and prepare a new validation dataset such as StanFi. Build the adapter, split manifests, and first baseline run path. \\
July 2026 & Run clean baseline and clean SupCon experiments on the new dataset. Compare whether the PA-CSI MultiEnv conclusions transfer under the same reporting discipline. \\
August 2026 & If the new dataset confirms a similar failure mode, design one representation/pseudo-label diagnostic and evaluate it on both PA-CSI MultiEnv and the new dataset. \\
September 2026 & Expand only the surviving method to additional seeds, LOEO or cross-subject targets, and geometry metrics. Add robustness checks only after cross-dataset behavior is understood. \\
October 2026 & Decide the final contribution boundary: cross-dataset method, diagnostic benchmark, or negative-result plus practical protocol paper. Draft the manuscript and supplement. \\
November 2026 & Finalize reproducibility materials, dataset adapters, run logs, split manifests, and submission package. Select venue path based on method strength and evaluation breadth. \\
\bottomrule
\end{longtable}

\subsection*{Target Venues}

\begin{longtable}{@{}p{1.25in}p{1.35in}p{3.55in}@{}}
\toprule
Venue & Priority & Reason \\
\midrule
\endhead
IMWUT / UbiComp & Primary long-term target & Best fit if the work becomes a strong sensing paper with convincing domain-generalization analysis, cross-environment evidence, and a clearly justified method. \\
Sensors & Practical fallback journal route & Suitable if the final contribution is a reproducible PA-CSI DG study with careful baselines, diagnostics, and moderate method novelty. \\
IEEE Access & Backup journal route & Useful if the main value is broad empirical evaluation, reproducibility, and engineering guidance rather than a highly novel algorithm. \\
\bottomrule
\end{longtable}

\section{Self-Assessment}

\paragraph{What went well.}
I made the research process stricter. May began with several plausible method ideas, but I did not keep them just because they had positive smoke results. The matched full checks, diagnostic reports, and abandonment memo make the current evidence much cleaner.

\paragraph{What still needs improvement.}
I still spent too much effort on method variations before fully answering the data/representation question. The next stage should start from the failure mode and only then choose the method. I also need to keep all comparisons matched by seed, config family, and selection rule from the beginning of each experiment.

\paragraph{Main priority going forward.}
The main priority is to resolve the $E2$ class $1/4$ collapse question while moving beyond a single dataset. If manual review indicates data or semantic ambiguity, the project should handle that explicitly. In parallel, June should introduce a new dataset such as StanFi to check whether the clean SupCon baseline and the May negative conclusions transfer. In both cases, the May stopped-method list should remain closed unless a new mechanism or cross-dataset evidence changes the premise.

\vfill
\begin{center}
End of Report
\end{center}

\end{document}
